\subsection{Infrared seam and sector bifurcation}

\begin{theorem}[Sector-dressed spectral scale bifurcation]
\label{thm:sector-dressed-spectral-scale-bifurcation}
The primitive response $\chiseed$ bifurcates on the closed branch into distinct dressed spectral
channels:
\[
\chiseed\longrightarrow \chigeozero
\qquad\text{and}\qquad
\chiseed\longrightarrow \chi_{\QCD}.
\]
Here $\chigeozero$ is the stationary Einstein/Higgs background scale, while $\chi_{\QCD}$ is the
transport-dressed hadronic scale. They share one boundary origin but are not identified on the
closed theorem surface.
\end{theorem}

\begin{proof}
The geometric channel is defined by the local spectral scale and its stationary branch
\Cref{def:local-spectral-scale,thm:stationary_geometric_scale_equation}. The hadronic channel is
kept separate by the sector-dressed transport/confinement bookkeeping recorded in the transport
block and summarized in the scale-bifurcation diagram \Cref{fig:chi-bifurcation-clean}. Hence the
same primitive response is dressed differently in the geometric and hadronic sectors.
\end{proof}

\begin{theorem}[Carrier-exhausted infrared seam eigenvalue under the exhaustion hypothesis]
\label{thm:carrier-exhausted-ir-seam-eigenvalue}
Let
\[
q_\star:=\deltatop e^{-2\alpha_\star},
\qquad
\phibase=\frac{1}{6\pi}.
\]
Assume the leading admissible infrared seam mode is rank one and exhausts the nonempty carrier
occupancy shadow, so that
\[
\rho(U_\Sigma)=\phibase\,q_\star^{2^{\gcar}-1}.
\]
Then on the canonical branch $\gcar=5$ and therefore
\[
\rho(U_\Sigma)=\phibase\,q_\star^{31}.
\]
Moreover
\[
\frac{\Lambda_{\mathrm{IR}}}{\Mplunred^4}
=
-\log\!\bigl(1-\rho(U_\Sigma)\bigr).
\]
\end{theorem}

\begin{proof}
The first displayed identity is exactly the carrier-exhaustion hypothesis written in closed-branch
variables. On the canonical carrier branch, \TFPTxref{cor:rank-5-master-compression} gives
$\gcar=5$, hence $2^{\gcar}-1=31$. The final identity is the rank-one case of
\TFPTxref{cor:ir-seam-transfer-bounds}.
\end{proof}

\subsection{Physical readout and metrology functor}

\begin{theorem}[Exact electroweak residue matching as a dimensionless readout]
\label{thm:dimensionless-ew-residue-matching}
On the renormalized branch, define
\[
\mathcal Z_{\mathrm{EW}}^{\mathrm{TFPT}}
:=
\frac{1}{2\vgeo^2\,\mathcal R_{\mathrm{CC}}^{\mathrm{TFPT}}(0)}
=
Z_\Phi(0)\left(1+\frac{\delta v}{\vgeo}\right)^2
\,\frac{1+\delta_W}{1+\Delta r_{\mathrm{CC}}^{\mathrm{TFPT}}}.
\]
Then the quantities
\[
\mathcal R_{\mathrm{CC}}^{\mathrm{TFPT}}(0),
\qquad
Z_\Phi(0),
\qquad
\delta v,
\qquad
\delta_W,
\qquad
\Delta r_{\mathrm{CC}}^{\mathrm{TFPT}}
\]
are determined by $\Gamma_{\mathrm{TFPT}}^{\mathrm{ren}}$. Consequently
\[
\frac{v_{\mathrm{phys}}}{\lambda_\Sigma}
=
\frac{\vgeo}{\lambda_\Sigma}\sqrt{\mathcal Z_{\mathrm{EW}}^{\mathrm{TFPT}}},
\qquad
\frac{v_{\mathrm{phys}}^2}{\lambda_\Sigma^2}
=
\left(\frac{\vgeo}{\lambda_\Sigma}\right)^2\mathcal Z_{\mathrm{EW}}^{\mathrm{TFPT}},
\]
and hence
\[
G_Nv_{\mathrm{phys}}^2
=
G_N\vgeo^2\,\mathcal Z_{\mathrm{EW}}^{\mathrm{TFPT}}
\]
is an internal readout of the closed branch.
\end{theorem}

\begin{proof}
The matching identity is exactly \TFPTxref{thm:ew-geometric-to-physical-matching} written in
boundary-normalized form. By \TFPTxref{thm:operational-completeness}, the same residue data belong to
the renormalized 1PI package $\Gamma_{\mathrm{TFPT}}^{\mathrm{ren}}$ before any scheme choice is
made. Multiplying by $G_N\lambda_\Sigma^2$ yields the final formula for $G_Nv_{\mathrm{phys}}^2$.
\end{proof}

\begin{corollary}[Dimensionless gravitational electroweak metrology]
\label{cor:dimensionless-gravitational-ew-metrology}
On the closed branch,
\[
G_Nv_{\mathrm{phys}}^2
=
\frac{1}{8\pi}\gcar^2\beta_{\mathrm{rad}}^4
\exp\!\left[-\frac{2(\alpha^{-1}(0)+\delta_{\mathrm{ph}})}{5}\right]
\,\mathcal Z_{\mathrm{EW}}^{\mathrm{TFPT}}.
\]
This is the primary gravity/electroweak hierarchy prediction of the closed branch. It contains no
external gravitational input and no metrological scale.
\end{corollary}

\begin{proof}
Combine \Cref{cor:dimensionless-geometric-hierarchy,thm:dimensionless-ew-residue-matching}.
\end{proof}

\begin{corollary}[Point-closed Planck-scale reconstruction]
\label{cor:point-closed-planck-reconstruction}
Using \Cref{thm:charged-current-point-closure}, the electroweak factor in
\Cref{cor:dimensionless-gravitational-ew-metrology} is fixed pointwise as
\[
\mathcal Z_{\mathrm{EW}}^{\mathrm{TFPT}}
=
1.0336068207261621.
\]
Consequently
\[
G_Nv_{\mathrm{phys}}^2
=
4.0671702198788166\times10^{-34}.
\]
After one non-gravitational metrological calibration of $v_{\mathrm{phys}}$, Newton's constant is
algebraically reconstructed by
\[
G_N=
\frac{(G_Nv_{\mathrm{phys}}^2)_{\mathrm{TFPT}}}{v_{\mathrm{phys}}^2},
\]
and the Planck units follow downstream:
\[
M_{\mathrm{Pl}}=\frac{1}{\sqrt{G_N}},
\qquad
\Mpl=\frac{1}{\sqrt{8\pi G_N}},
\qquad
\ell_P=\sqrt{G_N},
\qquad
t_P=\ell_P
\quad(c=\hbar=1).
\]
In SI representatives the usual conversion is applied only after this dimensionless hierarchy is
known:
\[
\ell_P=\sqrt{\frac{\hbar G_N}{c^3}},
\qquad
t_P=\sqrt{\frac{\hbar G_N}{c^5}},
\qquad
m_P=\sqrt{\frac{\hbar c}{G_N}}.
\]
\end{corollary}

\begin{proof}
The first displayed value is the point-closed charged-current readout of
\Cref{thm:charged-current-point-closure}. Multiplying it by the internally fixed geometric hierarchy
row $G_N\vgeo^2$ gives $G_Nv_{\mathrm{phys}}^2=G_N\vgeo^2\mathcal Z_{\mathrm{EW}}^{\mathrm{TFPT}}$.
Since this quantity is dimensionless, no Newton constant is inserted at this stage. Once a
non-gravitational representative for $v_{\mathrm{phys}}$ is chosen, the displayed algebraic formula
fixes $G_N$ and the Planck units follow by definition.
\end{proof}


\begin{theorem}[Pole masses as dimensionless branch readouts]
\label{thm:dimensionless-pole-mass-readout}
For each fermion sector $f$, the quantity
\[
\frac{m_f^{\mathrm{pole}}}{\lambda_\Sigma}
\]
is determined by the positive root of the Dirac pole equation of
\Cref{thm:dirac-pole-matching} after dividing by $\lambda_\Sigma^2$. Therefore
\[
\frac{m_f^{\mathrm{pole}}}{\lambda_\Sigma}
=
\mathcal M_f(\Tstar),
\qquad
G_N\bigl(m_f^{\mathrm{pole}}\bigr)^2
=
(G_N\lambda_\Sigma^2)\left(\frac{m_f^{\mathrm{pole}}}{\lambda_\Sigma}\right)^2.
\]
\end{theorem}

\begin{proof}
By \Cref{thm:pole-level-readout}, fermion pole masses are read from the dressed inverse propagators
of the renormalized branch and therefore belong to $\OphysTFPT$. Dividing the pole equation by
$\lambda_\Sigma^2$ produces a dimensionless root problem. The last identity is immediate.
\end{proof}

\begin{definition}[TFPT metrology functor]
\label{def:tfpt-metrology-functor}
Define the dimensionless metrology assignment of the closed branch by
\[
\mathfrak M_{\mathrm{TFPT}}(\Tstar)
:=
\left\{
\widehat O \;\middle|\;
O\in\OphysTFPT,\quad
d_O=\dim_{\mathrm{mass}}(O),\quad
\widehat O=\frac{O}{\lambda_\Sigma^{d_O}}
\right\}.
\]
\end{definition}

\begin{theorem}[Dimensionless TFPT metrology]
\label{thm:dimensionless-tfpt-metrology}
For every physical observable $O$ of mass dimension $d_O$ in the physical observable layer, one has
\[
\frac{O}{\lambda_\Sigma^{d_O}}
=
\mathcal F_O(\Tstar).
\]
In particular, the closed branch fixes
\[
\alpha,
\qquad
\delta_{\mathrm{ph}},
\qquad
\lambda_C,
\qquad
\theta_{\mathrm{eff}},
\qquad
G_N\lambda_\Sigma^2,
\qquad
\frac{v_{\mathrm{phys}}}{\lambda_\Sigma},
\qquad
G_Nv_{\mathrm{phys}}^2,
\]
\[
\frac{m_f^{\mathrm{pole}}}{\lambda_\Sigma},
\qquad
G_N\bigl(m_f^{\mathrm{pole}}\bigr)^2,
\qquad
\frac{\Lambda_{\mathrm{IR}}}{\lambda_\Sigma^4},
\qquad
\frac{\Lambda_{\mathrm{IR}}}{\Mplunred^4},
\qquad
\frac{m_i}{m_j}
\]
without external metrological input. Only afterwards may one choose a scheme representative or SI
translation.
\end{theorem}

\begin{proof}
By \Cref{thm:boundary-spectral-unit-rigidity}, the unit $\lambda_\Sigma$ is fixed by the one-sided
boundary datum. By \TFPTxref{thm:operational-completeness}, every main-text observable factors through
the physical observable functor
\[
\Mphys\circ\Rren:\mathbf{TFPT}^{\mathrm{rig}}\to\OphysTFPT.
\]
The dimensionless observables $\alpha$ and $\delta_{\mathrm{ph}}$ are already fixed directly by the
electromagnetic and transport closures. The quantities
$G_N\lambda_\Sigma^2$, $v_{\mathrm{phys}}/\lambda_\Sigma$, and
$G_Nv_{\mathrm{phys}}^2$ are fixed by
\Cref{cor:dimensionless-planck-metrology,thm:dimensionless-ew-residue-matching,cor:point-closed-planck-reconstruction}. The pole-mass
claims follow from \Cref{thm:dimensionless-pole-mass-readout}. The dimensionless cosmological
quantity $\Lambda_{\mathrm{IR}}/\Mplunred^4$ is fixed by the theorem-level determinant expression of
\TFPTxref{thm:cosmology-closure}, while mass ratios are already dimensionless quotients of pole
readouts. Therefore every displayed element is an internal function of $\Tstar$, and only the final
map to $\OschemeTFPT/\SchGrp$ chooses a comparison convention.
\end{proof}


% ---- Conditional minimal-parameter picture: source lines 13773-13838 ----
\subsection{Conditional minimal-parameter picture}

The next package is stronger but also more conditional. It assumes that the carrier rank, family multiplicity, and seam normalization are all read as closed rather than only partially inherited.

\begin{principle}[Conditional minimal-parameter compression]
Assume the canonical branch
\[
\gcar:=\rank E=5,
\qquad
N_{\mathrm{fam}}=3,
\qquad
\cthree=\frac{1}{8\pi}.
\]
Then the numerical kernel may be rewritten almost entirely in terms of
\[
(\pi,\gcar,N_{\mathrm{fam}})=(\pi,5,3).
\]
Indeed, on this branch one has
\[
\gamma=\frac56=\frac{\gcar}{\gcar+1},
\qquad
\dim S^+=2^{\gcar-1},
\qquad
\Oadm=N_{\mathrm{fam}}2^{\gcar-1}=48,
\]
\[
\phibase=\frac{1-\gamma}{\pi}=\frac{1}{(\gcar+1)\pi},
\qquad
\deltatop=\Oadm\,\cthree^4,
\qquad
\delta_2=\Dstart\,\Oadm\,\cthree^8,
\]
with
\[
\Dstart=\gcar\,\dim\mathfrak g_{\SM}=5\cdot 12=60,
\qquad
\Oadm=(\gcar-1)\dim\mathfrak g_{\SM}=4\cdot 12=48.
\]
\end{principle}

\begin{remark}[Compressed dependency chain]
Under the same derived closure package, the appendix-level numerical chain reads
\[
(\widetilde M\to M,\Seam,\taudbl,\iC,\chiseed)
\Longrightarrow
E_3\oplus E_2
\Longrightarrow
Y \leftrightarrow \epscar
\Longrightarrow
S(U(3)\times U(2))
\Longrightarrow
S^+,
\]
\[
\Longrightarrow
N_{\mathrm{fam}}=3
\Longrightarrow
\Oadm=48
\Longrightarrow
(\deltatop,\delta_2,\phiz)
\Longrightarrow
(\alpha,\lambda_C,\beta_{\mathrm{rad}},\Omega_b,\sin^2\theta_{13}),
\]
with the optional E$_8$, inflationary, and gravity-side sectors then read as downstream continuations of the same compressed core.
\end{remark}

